\begin{solapartat}
Sabem que l'objecte, paper, és a 10~cm de la lent, $s=-10\ \mathrm{cm}$.

\punts{0,25} Per trobar on es forma la imatge utilitzem l'equació de la lent: $\dfrac{1}{s'} - \dfrac{1}{s} = \dfrac{1}{f'}$

L'augment lateral és la relació entre l'alçada de la imatge i de l'objecte. $\beta = \dfrac{y'}{y}$

Igualment, com que els triangles són iguals també és la relació entre posicions de la imatge i l'objecte, $\beta = \dfrac{y'}{y} = \dfrac{s'}{s}$.

\punts{0,25} Ens diuen que l'augment lateral és 2: $\beta = \dfrac{y'}{y} = \dfrac{s'}{s} = 2$

Per tant, $s' = 2s = -20\ \mathrm{cm}$

\punts{0,25} Introduint els valors a l'equació de la lent: $\dfrac{1}{-20} - \dfrac{1}{-10} = \dfrac{1}{f'}$

\punts{0,25} Per tant $\dfrac{1}{-20} - \dfrac{1}{-10} = \dfrac{-10+20}{200} = \dfrac{10}{200} = \dfrac{1}{f'}$ la focal de la lent és $f' = 20\ \mathrm{cm}$. % DUBTE: a l'original el signe "+" de "-10+20" no es veu (probable error de tipografia del PDF)

\punts{0,25} La potència d'una lent (en diòptries) és la inversa de la focal (en metres). Per tant:
\[ P = \frac{1}{f'} = \frac{1}{0,2} = 5\ D. \]
\end{solapartat}

\begin{solapartat}
\punts{0,25} Hem de situar la lent a 22~cm de la taula per a que s'hi formi la imatge, per tant $s'=22\ \mathrm{cm}$.

\punts{0,25} Per trobar a quina distància es troba la làmpada de la taula, podem buscar quina és la posició de la làmpada respecte la lent utilitzant l'equació de la lent: $\dfrac{1}{s'} - \dfrac{1}{s} = \dfrac{1}{f'}$.

\punts{0,25} Introduint els valors a l'equació de la lent: $\dfrac{1}{22} - \dfrac{1}{s} = \dfrac{1}{20}$

\punts{0,25} Per tant $\dfrac{1}{-s} = \dfrac{1}{20} - \dfrac{1}{22} = \dfrac{22-20}{440} = \dfrac{2}{440}$, per tant $s=-220\ \mathrm{cm}$, que és la posició de la làmpada respecte la lent.

\punts{0,25} Per tant la distància de la làmpada a la taula és $|s' - s|$ o la suma de les dues distàncies: $|s'| + |s|$. Posant valors, $\textit{distància} = 220 + 22 = 242\ \mathrm{cm}$.
\end{solapartat}
